%% ====================================================================
\section{Numerical Methods}
\label{sec:methods}
%% ====================================================================

\subsection{Computation of Partial Dirichlet Series}
\label{sec:dirichlet}

The fundamental computational task is evaluating the partial Dirichlet series:
\begin{equation}\label{eq:partial-dirichlet}
  S_N(s) = \sum_{n=1}^{N} n^{-s}, \qquad s = \beta + it.
\end{equation}

\paragraph{Direct computation.}
Each term is computed as
\begin{equation}\label{eq:term-decomp}
  n^{-s} = n^{-\beta}\,n^{-it} = n^{-\beta}\,e^{-it\ln n},
\end{equation}
and separated into real and imaginary parts:
\begin{align}
  \re(n^{-s}) &= n^{-\beta}\cos(t\ln n), \label{eq:re-term}\\
  \im(n^{-s}) &= -n^{-\beta}\sin(t\ln n). \label{eq:im-term}
\end{align}
This decomposition avoids complex exponentiation and uses only real arithmetic (two
multiplications, one cosine, one sine per term), which is more numerically stable and
cache-friendly.

\paragraph{Alternating series.}
For System~1, we need the alternating partial sum
\begin{equation}\label{eq:eta-partial}
  \eta_N(s) = \sum_{n=1}^{N} (-1)^{n+1}\,n^{-s},
\end{equation}
obtained by multiplying each term by $(-1)^{n+1} = +1$ for odd~$n$, $-1$ for even~$n$.

\paragraph{Numerical stability.}
The terms $n^{-\beta}$ decrease monotonically for $\beta > 0$, ensuring numerical
stability of the running sum (no catastrophic cancellation in the absolute sense).
However, for the alternating series with small~$\beta$, consecutive terms nearly cancel,
leading to loss of relative precision. For $\beta = 0.5$ and large~$n$, consecutive
terms differ by $O(1/n)$, so the partial sum has relative precision degraded by a factor
of roughly $\sqrt{N}$. For $N = 10^5$, this costs approximately $2$--$3$ decimal digits,
which is acceptable in double precision ($\approx 15$--$16$ significant digits).

\paragraph{Precomputation.}
For a fixed~$\beta$, the values $n^{-\beta}$ ($n = 1, \ldots, N$) are independent of~$t$
and can be precomputed and stored. This reduces the per-point cost to $O(N)$
multiplications and trigonometric evaluations.

\paragraph{Error bound.}
The truncation error for the ordinary Dirichlet series is
\begin{equation}\label{eq:trunc-ordinary}
  |S_N(s) - \zeta(s)| \le \sum_{n=N+1}^{\infty} n^{-\beta}.
\end{equation}
By integral comparison:
\begin{equation}\label{eq:integral-bound}
  \sum_{n=N+1}^{\infty} n^{-\beta}
  \le \int_{N}^{\infty} x^{-\beta}\,dx
  = \frac{N^{1-\beta}}{\beta - 1}, \qquad \beta > 1.
\end{equation}
For $0 < \beta \le 1$ (the critical strip), the ordinary Dirichlet series does not
converge, but the alternating series converges. The truncation error for the alternating
series is bounded by the Dirichlet test (monotone coefficients, bounded partial sums of
$(-1)^{n+1}$):
\begin{equation}\label{eq:trunc-alt}
  |\eta_N(s) - \eta(s)| \le C(s)\,(N+1)^{-\beta},
\end{equation}
where $C(s) = O(1)$ for bounded $|s|$. For $\beta = 1/2$ and $N = 10^4$: error
$\lesssim 10^{-2}$; for $N = 10^6$: error $\lesssim 10^{-3}$. This is the fundamental
limitation of direct summation in the critical strip.

\paragraph{Euler--Maclaurin acceleration.}
For improved convergence, we apply the Euler--Maclaurin summation formula to the tail:
\begin{multline}\label{eq:euler-maclaurin}
  \sum_{n=N+1}^{M} n^{-s}
  \approx \int_{N}^{M} x^{-s}\,dx + \frac{1}{2}\bigl(N^{-s} + M^{-s}\bigr) \\
  + \sum_{k=1}^{p} \frac{B_{2k}}{(2k)!}\,
  s(s+1)\cdots(s+2k-2)\,\bigl(N^{-(s+2k-1)} - M^{-(s+2k-1)}\bigr) + R_p,
\end{multline}
where $B_{2k}$ are Bernoulli numbers and $R_p$ is the remainder. Taking $M \to \infty$
and using $p$~terms of the asymptotic expansion, the error bound must account for the
rising factorial $|s(s{+}1)\cdots(s{+}2p{-}2)|$, which grows as $|s|^{2p-1} \sim t^{2p-1}$
for large~$t$. The correct error estimate is
\begin{equation}\label{eq:em-error}
  |R_p| = O\bigl(|s|^{2p-1}\,N^{-(\beta+2p-1)}\bigr).
\end{equation}
This means that for large~$t$, the Euler--Maclaurin correction terms grow and the
asymptotic expansion eventually diverges (as is typical of asymptotic series). The
optimal number of correction terms $p$ depends on both $N$ and~$t$.

\paragraph{Precision table for Euler--Maclaurin correction ($\beta = 1/2$).}
In the following table, $|s|^{2p-1}$ denotes the magnitude of the rising factorial
$s(s{+}1)\cdots(s{+}2p{-}2)$, approximated as $|s|^{2p-1}$ for $|s| \gg 1$. With
$s = 1/2 + it$, we have $|s| \approx t$ for large~$t$.

\begin{table}[H]
\centering
\caption{Euler--Maclaurin effective error at $\beta = 1/2$.}
\label{tab:em-precision}
\begin{tabular}{@{}cccccc@{}}
\toprule
$N$ & $t$ & $p$ & $|s|^{2p-1}$ & $N^{-(\beta+2p-1)}$ & Effective error \\
\midrule
$10^4$ & $10^2$ & 3 & $10^{10}$  & $10^{-22}$ & $\sim 10^{-12}$ \\
$10^4$ & $10^2$ & 5 & $10^{18}$  & $10^{-38}$ & $\sim 10^{-20}$ \\
$10^4$ & $10^4$ & 3 & $10^{20}$  & $10^{-22}$ & $\sim 10^{-2}$  \\
$10^4$ & $10^4$ & 5 & $10^{36}$  & $10^{-38}$ & $\sim 10^{-2}$  \\
$10^5$ & $10^2$ & 5 & $10^{18}$  & $10^{-47.5}$ & $\sim 10^{-29.5}$ \\
$10^3$ & $10^2$ & 3 & $10^{10}$  & $10^{-16.5}$ & $\sim 10^{-6.5}$ \\
\bottomrule
\end{tabular}
\end{table}

The table reveals that for large~$t$, the Euler--Maclaurin correction becomes unreliable:
at $t = 10^4$ with any~$p$, the rising factorial overwhelms the $N^{-(\beta+2p-1)}$ decay.
This is the fundamental reason why the Riemann--Siegel formula (Section~\ref{sec:rs}) is
essential for $t > 10^3$. For moderate~$t$ ($t < 10^3$), Euler--Maclaurin with $p = 5$
and $N = 10^4$ provides $\sim 20$ digits of precision, far exceeding double-precision
limits.

\begin{remark}[On the relationship between acceleration and the DQPT framework]
\label{rem:acceleration-circularity}
A methodological subtlety deserves comment. The DQPT framework is formulated in terms of
the thermodynamic limit $N \to \infty$ of the \emph{finite} partial sums
$\mathcal{L}_N(\beta,t)$, and the physics resides in the finite-$N$ scaling behavior
(Section~\ref{sec:system1}). When we apply Euler--Maclaurin summation or Borwein
acceleration to compute the ``true'' value $\eta(s)$ (the infinite-sum limit), we are
effectively bypassing the finite-$N$ physics and returning to standard analytic number
theory. This does not invalidate the computation, but it means that the accelerated
results should be understood as \emph{high-precision evaluations of the same mathematical
function}, not as simulations of a physical system at finite~$N$. We use both approaches:
unaccelerated partial sums (to verify finite-$N$ scaling predictions of the DQPT theory)
and accelerated evaluations (to achieve the precision needed for reliable zero detection
and comparison with published tables).
\end{remark}

\paragraph{Borwein--Cohen--Villegas--Zagier acceleration.}
For the eta function specifically, Borwein~\cite{borwein2000} showed that the following
identity provides geometric convergence:
\begin{equation}\label{eq:borwein}
  \eta(s) = -\frac{1}{d_n}\sum_{k=0}^{n-1}(-1)^k\,(d_k - d_n)\,(k+1)^{-s},
\end{equation}
where the coefficients $d_k$ are defined by
\begin{equation}\label{eq:borwein-dk}
  d_k = n\sum_{j=0}^{k} \frac{(n+j-1)!\,4^j}{(n-j)!\,(2j)!}
\end{equation}
and $d_n = d_k|_{k=n}$.

\paragraph{Numerical stability of the coefficients.}
For $n = 40$, the coefficients $d_k$ grow extremely rapidly ($d_{40} \sim 10^{23}$),
exceeding the range of IEEE~754 double precision (which represents integers exactly only
up to $2^{53} \approx 9 \times 10^{15}$). To avoid overflow and catastrophic
cancellation, we compute the \emph{ratios} $c_k = (d_k - d_n)/d_n$ directly via a
stable recurrence. Define $e_k = d_k/d_n$; then $c_k = e_k - 1$, and the Borwein
formula becomes
\begin{equation}\label{eq:borwein-stable}
  \eta(s) = -\sum_{k=0}^{n-1}(-1)^k\,c_k\,(k+1)^{-s}.
\end{equation}
The ratios $e_k$ can be computed via the recurrence relation for the $d_k$ sequence,
working entirely in double precision without overflow (since $0 \le e_k \le 1$ for all
$k < n$).

\paragraph{Error bound and region of validity.}
The approximation error for the eta function satisfies~\cite{borwein2000}
\begin{equation}\label{eq:borwein-error}
  |\eta(s) - \text{approximation}| \le 3\,(3 + \sqrt{8}\,)^{-n}
\end{equation}
for $s$ in the region $\re(s) > 0$. With $n = 40$:
\begin{equation}\label{eq:borwein-40}
  |\text{error}_\eta| < 3 \times 5.828^{-40} \approx 10^{-30}.
\end{equation}
If one needs $\zeta(s)$ instead, it is recovered via $\zeta(s) = \eta(s)/(1 - 2^{1-s})$.
The resulting error in $\zeta$ is amplified by the factor $1/|1 - 2^{1-s}|$, which
diverges near $s = 1 + 2\pi ik/\log 2$ for $k \in \Z$. For computations in the critical
strip ($0 < \re(s) < 1$), $|1 - 2^{1-s}|$ is bounded away from zero. To see this:
$2^{1-s} = 1$ requires both $|2^{1-s}| = 2^{\re(1-s)} = 1$ (i.e., $\re(1-s) = 0$,
hence $\re(s) = 1$) and $\arg(2^{1-s}) = \im(1-s)\log 2 \in 2\pi\Z$. Since $\re(s) = 1$
lies outside the open critical strip $0 < \re(s) < 1$, the prefactor $1 - 2^{1-s}$ never
vanishes inside the strip. Moreover, for $\re(s) \in [0.1, 0.9]$,
$|1 - 2^{1-s}| \ge c > 0$ for some explicit constant~$c$ (computable for each~$\beta$),
so the error amplification is bounded.

This geometric convergence rate of $O(5.828^{-n})$ means $n = 40$ terms suffice for
$\sim 30$ digits of accuracy---far exceeding double precision requirements. This is our
primary method for computing $\eta(s)$ in the critical strip.

\paragraph{Kahan compensated summation.}
All summation loops use the Kahan compensated summation algorithm~\cite{kahan1965} to
reduce the accumulated floating-point rounding error from $O(N\,\eps_{\mathrm{mach}})$
to $O(\eps_{\mathrm{mach}})$, independent of~$N$. The algorithm maintains a running
compensation variable that tracks the low-order bits lost in each addition:
\begin{lstlisting}[language=C++,caption={Kahan compensated summation.}]
double sum = 0.0, c = 0.0;
for (int i = 0; i < N; ++i) {
    double y = x[i] - c;
    double t = sum + y;
    c = (t - sum) - y;
    sum = t;
}
\end{lstlisting}

\begin{remark}[FMA and GPU execution]\label{rem:fma}
CUDA compiles with fused multiply-add (FMA) instructions enabled by default. FMA
computes $a \cdot b + c$ with a single rounding (rather than the two roundings of
separate multiply and add), which changes the numerical result at the level of
$O(\eps)$ per operation. This has two consequences:
\begin{enumerate}
  \item The Kahan summation algorithm relies on the rounding behavior of individual
        additions. With FMA enabled, the compiler may fuse the compensation steps,
        potentially defeating the purpose of the algorithm. To ensure correct Kahan
        summation on GPU, use \texttt{\_\_dadd\_rn()} and \texttt{\_\_dsub\_rn()}
        intrinsics (round-to-nearest double-precision add/subtract) or compile with
        \texttt{--fmad=false} for the relevant kernels.
  \item CPU and GPU results will differ at $O(\eps)$ per operation due to different FMA
        behavior (CPU code compiled with \texttt{-ffp-contract=off} does not use FMA).
        When comparing CPU and GPU results for validation, the expected discrepancy is
        $O(N\,\eps)$ without Kahan summation, or $O(\eps)$ with correct Kahan summation
        on both platforms.
\end{enumerate}
\end{remark}


\subsection{Riemann--Siegel Theta Function}
\label{sec:theta}

The Riemann--Siegel theta function is defined by
\begin{equation}\label{eq:theta-def}
  \vartheta(t) = \im\!\bigl(\log\Gamma\bigl(\tfrac{1}{4} + \tfrac{it}{2}\bigr)\bigr)
  - \frac{t}{2}\log\pi.
\end{equation}

\paragraph{Stirling asymptotic expansion.}
For $|t| \gg 1$, we use the asymptotic expansion of $\log\Gamma$ via Stirling's series.
Setting $z = 1/4 + it/2$:
\begin{equation}\label{eq:stirling}
  \log\Gamma(z) = \Bigl(z - \tfrac{1}{2}\Bigr)\log z - z + \tfrac{1}{2}\log(2\pi)
  + \sum_{k=1}^{p} \frac{B_{2k}}{2k(2k-1)\,z^{2k-1}} + R_p(z).
\end{equation}
Taking the imaginary part and substituting $z = 1/4 + it/2$, we
obtain~\cite{edwards1974}\footnote{The first two correction terms ($1/(48t)$ and
$7/(5760t^3)$) are well known; the higher-order coefficients follow from the Stirling
series and are tabulated in Edwards~\cite{edwards1974}, Section~6.5.}
\begin{equation}\label{eq:theta-asymp}
  \vartheta(t) = \frac{t}{2}\ln\frac{t}{2\pi} - \frac{t}{2} - \frac{\pi}{8}
  + \frac{1}{48t} + \frac{7}{5760t^3} + \frac{31}{80640t^5}
  - \frac{127}{430080t^7} + O(t^{-9}).
\end{equation}
For our implementation, we use the first four correction terms (through $t^{-7}$), which
provides more than adequate precision for $t \ge 10$.

\paragraph{Error bound.}
The error after $p$ terms of the Stirling expansion is bounded by
\begin{equation}\label{eq:stirling-error}
  |R_p| \le \frac{|B_{2p+2}|}{(2p+2)(2p+1)\,|z|^{2p+1}}.
\end{equation}
For $t > 10$ and $p = 4$ (using terms through $t^{-7}$): $|R_4| < 10^{-12}$, which is
adequate for double precision. For $t > 50$ and $p = 3$ (through $t^{-5}$):
$|R_3| < 10^{-14}$.

\paragraph{Validity range.}
The asymptotic expansion is valid for $t > 0$ but achieves useful precision only for
$t > t_{\min}$. We require:
\begin{itemize}
  \item For double precision ($10^{-15}$): $t_{\min} \approx 10$ (with 4 correction terms).
  \item For extended precision ($10^{-30}$): $t_{\min} \approx 5$ (with 8 correction terms).
\end{itemize}
For $t < t_{\min}$, we fall back to numerical evaluation of $\log\Gamma$ using the
Lanczos approximation or a rational Chebyshev approximation for~$\Gamma$ in the relevant
region.

\begin{remark}[Branch cuts and small~$t$]\label{rem:branch}
The function $\vartheta(t) = \im(\log\Gamma(1/4 + it/2)) - (t/2)\log\pi$ involves the
complex logarithm of the Gamma function, which has branch cuts. For small~$t$ ($t < 10$),
the argument $1/4 + it/2$ is close to the real axis where $\Gamma$ has poles at
$0, -1, -2, \ldots$, and the principal branch of $\log\Gamma$ requires careful handling
to avoid discontinuities. Our implementation strategy:
\begin{itemize}
  \item For $t \ge 10$: use the Stirling asymptotic expansion ($4+$ correction terms),
        which is well away from branch cut issues and provides adequate precision.
  \item For $t < 10$: use the Lanczos approximation for $\Gamma(z)$ (with reflection
        formula for $\re(z) < 0.5$), or equivalently use the recurrence
        $\Gamma(z{+}1) = z\,\Gamma(z)$ to shift the argument into a region where the
        Lanczos series converges rapidly. This avoids the branch cut issues of
        $\log\Gamma$ by computing $\Gamma$ directly and then taking the logarithm.
\end{itemize}
In practice, the first zeta zero is at $t \approx 14.13$, so the small-$t$ regime is
only relevant for theta function evaluation at sub-zero-height positions (used in Gram
point computation and the Riemann--von Mangoldt formula).
\end{remark}


\subsubsection{The Riemann--Siegel Formula}
\label{sec:rs}

The Riemann--Siegel (RS) formula provides an efficient method for computing the Hardy
$Z$-function $Z_H(t) = e^{i\vartheta(t)}\zeta(1/2 + it)$ directly, requiring only
$O(\sqrt{t/(2\pi)})$ terms rather than $O(t)$ terms for direct
summation~\cite{titchmarsh1986,gabcke1979}. This is the standard method for large-scale
zero computations.

\paragraph{Explicit formula.}
The RS formula is
\begin{equation}\label{eq:rs-formula}
  Z_H(t) = 2\sum_{n=1}^{M} n^{-1/2}\cos\bigl(\vartheta(t) - t\ln n\bigr) + R(t),
\end{equation}
where $M = \lfloor\sqrt{t/(2\pi)}\rfloor$ is the main sum cutoff and $R(t)$ is the
remainder term.

\paragraph{First correction term.}
The leading-order remainder is given by the Riemann--Siegel
correction~\cite{gabcke1979}:
\begin{equation}\label{eq:rs-R0}
  R_0(t) = (-1)^{M-1}\Bigl(\frac{t}{2\pi}\Bigr)^{-1/4}\Psi(p),
\end{equation}
where $p = \sqrt{t/(2\pi)} - M$ is the fractional part and $\Psi(p)$ is the
Riemann--Siegel function:
\begin{equation}\label{eq:rs-psi}
  \Psi(p) = \frac{\cos\bigl(2\pi(p^2 - p - 1/16)\bigr)}{\cos(2\pi p)}.
\end{equation}
Higher-order correction terms $R_1, R_2, \ldots$ involve derivatives of~$\Psi$ and
powers of $(t/(2\pi))^{-1/2}$, each successive term reducing the error by an additional
factor of $O(t^{-1/2})$:
\begin{equation}\label{eq:rs-expansion}
  R(t) = R_0(t) + R_1(t) + R_2(t) + \cdots,
\end{equation}
where $|R_k(t)| = O(t^{-(2k+3)/4})$ (Siegel, 1932; Gabcke~\cite{gabcke1979}).

\paragraph{Error bounds after successive corrections.}
The following bounds are verified against
Titchmarsh~\cite{titchmarsh1986} and Gabcke~\cite{gabcke1979}:
\begin{align}
  \text{Main sum only:}\qquad &|Z_H(t) - \text{main sum}| = O(t^{-3/4}),
    \label{eq:rs-err0}\\
  \text{Main sum} + R_0:\qquad &|\text{error}| = O(t^{-5/4}),
    \label{eq:rs-err1}\\
  \text{Main sum} + R_0 + R_1:\qquad &|\text{error}| = O(t^{-7/4}).
    \label{eq:rs-err2}
\end{align}
Quantitatively (Gabcke~\cite{gabcke1979}): the error of the main sum alone is less than
$(3/2)\,t^{-3/4}$ for $t > 125$. Including the $R_0$ correction, the error is less than
$0.053\,t^{-5/4}$ for $t \ge 200$.

For $t = 10^4$: the uncorrected remainder is $O(t^{-3/4}) = O(10^{-3})$. Including
$R_0$, the error improves to $O(t^{-5/4}) = O(10^{-5})$. Including $R_0 + R_1$, the
error is $O(t^{-7/4}) = O(10^{-7})$.

\paragraph{Computational advantage.}
The main sum requires only $M = \lfloor\sqrt{t/(2\pi)}\rfloor$ terms:
\begin{center}
\begin{tabular}{@{}ll@{}}
\toprule
$t$ & $M$ \\
\midrule
$10^4$ & $\sim 40$ \\
$10^6$ & $\sim 399$ \\
$10^8$ & $\sim 3989$ \\
$10^{10}$ & $\sim 39894$ \\
\bottomrule
\end{tabular}
\end{center}
Compare this to direct summation of the Dirichlet series, which requires $N \gg t$ for
adequate convergence. The RS formula is essential for computations at large height.

\paragraph{Implementation notes.}
Each term in the main sum requires evaluation of
$n^{-1/2}\cos(\vartheta(t) - t\ln n)$. Since $n^{-1/2}$ is independent of~$t$, it can
be precomputed. The argument $\vartheta(t) - t\ln n$ must be reduced modulo $2\pi$
carefully to avoid loss of precision for large~$t$ (see Section~\ref{sec:trig} on
trigonometric evaluation).

\paragraph{Small-$t$ regime and method selection.}
The RS formula is an asymptotic approximation that improves with increasing~$t$.
At small heights, the main sum cutoff $M = \lfloor\sqrt{t/(2\pi)}\rfloor$ yields
very few terms:
\begin{center}
\begin{tabular}{@{}rrl@{}}
\toprule
$t$ & $M$ & RS viability \\
\midrule
$14$ & $1$ & unusable (single term) \\
$50$ & $2$ & poor \\
$100$ & $3$ & marginal \\
$200$ & $5$ & adequate (Gabcke bound applies) \\
$1000$ & $12$ & good \\
\bottomrule
\end{tabular}
\end{center}
For $t < 200$ (where the quantitative Gabcke error bound $0.053\,t^{-5/4}$ does not
apply), we compute the Hardy $Z$-function directly via
\begin{equation}\label{eq:zh-direct}
  Z_H(t) = \re\bigl[e^{i\vartheta(t)}\,\zeta\bigl(\tfrac{1}{2}+it\bigr)\bigr],
\end{equation}
where $\zeta(1/2+it)$ is obtained from the Borwein-accelerated eta function
(Section~\ref{sec:dirichlet}) as $\zeta(s) = \eta(s)/(1-2^{1-s})$. This provides
${\sim}30$ digits of precision independent of~$t$ (Paper Eq.~\ref{eq:borwein-40}).
For $t \ge 200$, the RS formula~\eqref{eq:rs-formula} with $R_0$ correction is used,
as it requires only $O(\sqrt{t})$ operations compared to the fixed cost of Borwein
($n = 100$ terms).

This two-method approach mirrors the strategy of Wei et~al.~\cite{wei2025}, who use
direct quantum simulation at small~$t$ and note that ``the deviation between estimated
and exact zeros decreases with increasing~$t$'' due to the improving RS approximation.
In our classical implementation, the Borwein method replaces the direct quantum
measurement for small~$t$.


\subsection{Zero Detection}
\label{sec:zero-detection}

\subsubsection{Bisection on the Hardy $Z$-function}
\label{sec:bisection}

The Hardy $Z$-function $Z_H(t)$ is real-valued, and its sign changes correspond to
simple zeros of $\zeta(1/2 + it)$. Our primary zero-detection method is:
\begin{enumerate}
  \item Evaluate $Z_H(t)$ on a uniform grid $t_0, t_0 + \Delta t, t_0 + 2\Delta t, \ldots$
  \item Identify intervals $[t_k, t_{k+1}]$ where $Z_H$ changes sign.
  \item Refine each sign change by bisection to precision~$\eps$.
\end{enumerate}

\paragraph{Grid spacing.}
The average spacing between consecutive zeros near height~$T$ is
\begin{equation}\label{eq:zero-spacing}
  \delta_{\mathrm{zero}} \sim \frac{2\pi}{\log(T/(2\pi))}.
\end{equation}
For $T = 100$: $\delta_{\mathrm{zero}} \approx 2.3$. For $T = 1000$:
$\delta_{\mathrm{zero}} \approx 1.2$. For $T = 10000$:
$\delta_{\mathrm{zero}} \approx 0.86$. We choose $\Delta t = \delta_{\mathrm{zero}}/4$
to ensure we do not miss any zeros (assuming zeros are roughly uniformly spaced, which
is a reasonable heuristic for the initial scan).

\subsubsection{Gram's Law and Its Failures}
\label{sec:gram}

Gram points are the values $g_n$ defined by $\vartheta(g_n) = n\pi$ for
$n = 0, 1, 2, \ldots$ Gram's ``law'' (actually a heuristic) states that
$(-1)^n Z_H(g_n) > 0$, i.e., $Z_H$ changes sign between consecutive Gram
points. This holds for approximately $73\%$ of Gram
intervals~\cite{trudgian2011}.

We use Gram points as an initial grid but do not rely on Gram's law for correctness.
Instead, we use finer subdivision wherever Gram's law fails (i.e., where $Z_H$ has the
same sign at two consecutive Gram points).

\subsubsection{Turing's Method for Zero Count Certification}
\label{sec:turing}

Turing~\cite{turing1953} developed a method to certify that all zeros up to height~$T$
have been found. The method relies on the Riemann--von Mangoldt formula for the
zero-counting function:
\begin{equation}\label{eq:rvm}
  N(T) = \frac{T}{2\pi}\log\frac{T}{2\pi} - \frac{T}{2\pi} + \frac{7}{8}
  + S(T) + O\!\left(\frac{1}{T}\right),
\end{equation}
where $S(T) = (1/\pi)\arg\zeta(1/2 + iT)$ satisfies $|S(T)| = O(\log T)$
unconditionally\footnote{See Trudgian~\cite{trudgian2011} for modern explicit bounds.
Note that $S(T)$ is \emph{not} bounded by a constant; the bound $|S(T)| = O(\log T)$
reflects the irregular distribution of zeta zeros.}. If the number of detected sign
changes of $Z_H$ in $[0, T]$ agrees with the nearest integer to
$(T/(2\pi))\log(T/(2\pi)) - T/(2\pi) + 7/8$, then all zeros in that range have been
found (up to the uncertainty introduced by $S(T)$).

We implement this check as a post-hoc validation: after detecting zeros up to
height~$T$, we verify that the count matches the Riemann--von Mangoldt prediction.

\subsubsection{Expected Zero Density}
\label{sec:zero-density}

The number of nontrivial zeros with $0 < \im(s) \le T$ is given by the Riemann--von
Mangoldt formula~\eqref{eq:rvm}. The leading terms give the approximation:
\begin{equation}\label{eq:rvm-approx}
  N(T) \approx \frac{T}{2\pi}\log\frac{T}{2\pi} - \frac{T}{2\pi} + \frac{7}{8}.
\end{equation}
This gives:
\begin{center}
\begin{tabular}{@{}rl@{}}
\toprule
$T$ & $N(T)$ \\
\midrule
$100$ & $\sim 29$ \\
$1{,}000$ & $\sim 649$ \\
$10{,}000$ & $\sim 10{,}142$ \\
$100{,}000$ & $\sim 138{,}069$ \\
\bottomrule
\end{tabular}
\end{center}


\subsection{Convergence Analysis}
\label{sec:convergence}

\subsubsection{Convergence of $L(\beta, t)$ with $N$}
\label{sec:conv-L}

For fixed $s = \beta + it$ with $0 < \beta < 1$, the alternating partial sum
$\eta_N(s)$ converges to $\eta(s)$ as $N \to \infty$. The rate of convergence is
\begin{equation}\label{eq:conv-rate}
  |\eta_N(s) - \eta(s)| = O(N^{-\beta}).
\end{equation}
This is the dominant source of error. For $\beta = 1/2$:
\begin{center}
\begin{tabular}{@{}rl@{}}
\toprule
$N$ & Error \\
\midrule
$10^4$ & $\sim 10^{-2}$ \\
$10^6$ & $\sim 10^{-3}$ \\
$10^8$ & $\sim 10^{-4}$ \\
\bottomrule
\end{tabular}
\end{center}
This slow convergence is the fundamental computational bottleneck. To achieve precision
$\delta$, we need $N \sim \delta^{-1/\beta} = \delta^{-2}$ for $\beta = 1/2$. This is
much worse than the Riemann--Siegel formula, which achieves precision~$\delta$ with
$O(\sqrt{T})$ terms, but the direct approach has the virtue of simplicity and direct
connection to the physical (DQPT) interpretation.

\paragraph{Improvement via Euler--Maclaurin.}
With $p$ terms of Euler--Maclaurin correction applied to the tail sum beyond~$N$, the
effective error becomes $O(|s|^{2p-1}\,N^{-(\beta+2p-1)})$ (see
Section~\ref{sec:dirichlet} for the full bound and precision table). For $p = 5$,
$\beta = 1/2$, and moderate~$t$ ($t < 10^4$): the error is well below double-precision
limits for $N \ge 10^4$. This makes Euler--Maclaurin-corrected summation practical for
moderate~$t$. For large~$t$, the Riemann--Siegel formula (Section~\ref{sec:rs}) is
preferable.

\subsubsection{Required $N$ as a Function of $t$ and Precision}
\label{sec:required-N}

The terms $n^{-s} = n^{-\beta}e^{-it\ln n}$ oscillate with frequency proportional
to~$t$. The Nyquist criterion requires that we sample these oscillations adequately, but
since we are summing (not sampling) the series, the oscillation frequency does not
directly increase the required~$N$. However, near zeros of~$\zeta$, the partial sums
exhibit cancellation at scale $O(N^{-\beta})$, so detecting zeros to relative
precision~$\delta$ requires
\begin{equation}\label{eq:required-N}
  N \ge \delta^{-1/\beta}.
\end{equation}
For $\beta = 1/2$ and $\delta = 10^{-6}$ (sufficient to distinguish a zero from a
near-miss): $N \ge 10^{12}$. This is impractical without acceleration techniques.

With Euler--Maclaurin ($p = 5$) and accounting for the $|s|^{2p-1}$ factor: the required
$N$ satisfies $|s|^{2p-1}\,N^{-(\beta+2p-1)} < \delta$, giving
\begin{equation}\label{eq:required-N-em}
  N \ge \Bigl(\frac{|s|^{2p-1}}{\delta}\Bigr)^{1/(\beta+2p-1)}.
\end{equation}
For $t = 100$, $|s| \approx 100$, $p = 5$, $\beta = 0.5$:
$N \ge (10^{18}/10^{-6})^{1/9.5} = 10^{24/9.5} \approx 336$. For $t = 10^4$:
$N \ge (10^{36}/10^{-6})^{1/9.5} = 10^{42/9.5} \approx 26{,}300$. These are achievable,
showing that Euler--Maclaurin correction is essential for practical computation at high
precision, though the required $N$ grows with~$t$.

\subsubsection{Computational Complexity}
\label{sec:complexity}

Direct summation: $O(N)$ operations per $(\beta, t)$ point. For $M$ grid points in the
$(\beta, t)$ plane, the total cost is $O(N \cdot M)$.

With Euler--Maclaurin: $O(N + p)$ per point, where $p$ is small ($5$--$10$). For a
moderate $N$ ($10^3$ to $10^4$), the cost is dominated by the direct sum, giving
$O(N \cdot M)$ total.

For a grid of $M_\beta = 100$ values of~$\beta$ and $M_t = 10{,}000$ values of~$t$,
with $N = 10^4$:
\begin{equation}\label{eq:total-ops}
  \text{Total operations} \sim 10^4 \times 10^6 = 10^{10}.
\end{equation}
On a single CPU core, this workload requires substantial time (minutes to hours 
depending on architecture and instruction set). While GPU parallelism 
drastically reduces wall-clock time, the hardware-level bottleneck of 
FP64 transcendental units must be considered (Section~\ref{sec:throughput}).


\subsection{GPU Parallelization Strategy}
\label{sec:gpu}

\subsubsection{Parallelism Structure}
\label{sec:parallel}

The computation of $L(\beta, t)$ or $G(\beta, t)$ at each $(\beta, t)$ point is
completely independent of all other points. This is an embarrassingly parallel workload,
ideal for GPU execution.

\paragraph{Parallelization scheme.}
\begin{itemize}
  \item Each CUDA thread computes one $(\beta, t)$ point.
  \item The thread iterates over $n = 1, \ldots, N$, accumulating the partial sum.
  \item The result (a complex number for~$L$, a real number for~$G$) is written to global
        memory.
\end{itemize}

\subsubsection{Memory Model}
\label{sec:memory}

For each $\beta$ value, we precompute the array $\{n^{-\beta}\}_{n=1}^{N}$ and store it
in global memory (or texture memory for cached access). Each thread reads this array and
computes the trigonometric factors $\cos(t\ln n)$ and $\sin(t\ln n)$ on the fly.

\paragraph{Memory requirements.}
\begin{itemize}
  \item Precomputed $n^{-\beta}$: $N \times \texttt{sizeof(double)} = N \times 8$ bytes
        per $\beta$ value. For $N = 10^5$ and $100$~$\beta$ values: $80$~MB (fits in
        $6$~GB VRAM).
  \item Precomputed $\ln(n)$: $N \times 8$ bytes $= 0.8$~MB (shared across all $\beta$
        values).
  \item Output array: $M_\beta \times M_t \times \texttt{sizeof(double)}$ per output
        quantity. For $100 \times 10{,}000 = 10^6$ points: $8$~MB.
  \item Total VRAM usage: $< 100$~MB, well within $6$~GB.
\end{itemize}

\subsubsection{Expected Throughput}
\label{sec:throughput}

RTX~2060 specifications:
\begin{itemize}
  \item 1920 CUDA cores.
  \item 6~GB GDDR6 VRAM.
  \item Base clock: 1365~MHz.
  \item FP64 throughput: approximately $210$~GFLOPS ($1/32$ of FP32).
\end{itemize}

Since our computation involves double-precision transcendental functions (\texttt{cos},
\texttt{sin}, \texttt{exp}, \texttt{log}), the throughput is limited by the special
function units (SFU).

\begin{remark}[SFU throughput for transcendentals]\label{rem:sfu}
On NVIDIA GPUs, the SFU throughput for double-precision transcendental functions is
approximately $1/4$ of the FP64 ALU throughput (and for some transcendentals, up to
$1/10$ due to multi-pass evaluation). Since our workload is transcendental-heavy (each
term requires one \texttt{cos} and one \texttt{sin} evaluation), the effective throughput
must be derated by $4$--$10\times$ compared to the peak FP64 rate.
\end{remark}

\paragraph{Revised throughput estimates.}
\begin{itemize}
  \item Per-point cost: $N$ multiplications $+ N$ transcendental evaluations. The
        transcendentals dominate: effective cost $\sim 10N$ FLOPS equivalent (accounting
        for SFU pipeline occupancy).
  \item For $N = 10^4$: $\sim 10^5$ FLOPS-equivalent per point.
  \item Effective throughput:
        $(210 \times 10^9/4)/(10^5) \approx 5 \times 10^5$ points/second (conservative;
        actual throughput depends on occupancy and memory access patterns).
\end{itemize}

For a $180 \times 10{,}000$ grid ($1.8 \times 10^6$ points): approximately $3$--$4$
seconds. For a $180 \times 100{,}000$ grid ($1.8 \times 10^7$ points): approximately
$30$--$40$ seconds. For $N = 10^5$ and a $180 \times 10{,}000$ grid: approximately
$30$--$40$ seconds.

These revised estimates ($4$--$10\times$ slower than naive FP64 peak estimates) still show
that GPU-accelerated computation is more than sufficient for our parameter ranges. Even
with additional overheads (memory latency, register pressure, thread divergence), runtimes
remain under a few minutes for the full sweep.

\subsubsection{Batch Size Estimation}
\label{sec:batch}

We organize the computation into batches to manage VRAM and kernel launch overhead:
\begin{itemize}
  \item Batch size: $10^5$ to $10^6$ $(\beta, t)$ points per kernel launch.
  \item Number of batches: total grid size / batch size.
  \item For a $100 \times 100{,}000$ grid: $10^7 / 10^6 = 10$ batches.
\end{itemize}
Each batch completes in $< 1$~second, so the total wall-clock time (including data
transfer) is under $10$~seconds for the full sweep.
